\section{A universal algebraic certificate}\label{sec:certificate}\label{cert:section}

The comparison needed for the stability inequality is a statement about
finite-dimensional tensors. We state it with free algebraic data before
deriving its geometric consequences. All tensor spaces here are over $\R^5$,
with the Euclidean inner product and full tensor norm; repeated indices range
from $1$ to $5$ and are summed.  A symmetric tensor is called trace-free if
contraction of any two
slots is zero.  We write $\STF^r(\R^5)$ for the space of symmetric trace-free
tensors of rank $r$.  Fix
\begin{equation}\label{cert:constants}
 a=\frac{29}{50},\qquad Q=250000,\qquad
 \epsilon_* =\frac{438907}{12000000},\qquad
 \delta_* =\frac{6984137359}{31250000000}.
\end{equation}

Let $A$ be a symmetric trace-free matrix with $|A|^2=\tr A^2=1$, and let
$X$ be a symmetric trace-free three-tensor.  Set
\[
 q_l=X_{lij}A_{ij},\qquad \chi=1-|p|^2-g^2,
\]
where $p,\ell\in\R^5$ and $g,w,z\in\R$ are otherwise free.  In particular,
$w$ and $z$ need not be positive.  On tensor polynomials in $A,p,g,\delta$,
define $D_l$ by the product rule and
\begin{equation}\label{cert:D}
 D_l A_{ij}=X_{lij}-q_lA_{ij},\qquad
 D_l p_i=wA_{li}-gz\delta_{li},\qquad
 D_l g=zp_l,\qquad D_l\delta_{ij}=0.
\end{equation}
This is a first derivative with values in the enlarged polynomial algebra
containing $X,w,z$; no derivative of these new variables is used.

\begin{proposition}[Universal inequalities]\label{cert:universal}
There are rational tensor polynomials $f,U_{jk},K_{lijk},M_l,P_l$ in
$A,p,g,\delta$, with the following properties. They are equivariant under
orthogonal changes of coordinates, with $f$ invariant. Their coefficients
are specified by Table~I in Appendix~\ref{app:tables}, using the decoding in
Section~\ref{cert:decode}. Put
\[
 C_{lijk}=K_{lijk}-K_{iljk},\qquad
 G_{lijk}=\delta_{li}U_{jk}+C_{lijk}.
\]
For every choice of the data above with $\chi\ge0$,
\begin{align}
 \mathcal I:={}&z^2\left[5U_{jk}A_{jk}
       +C_{lijk}(\delta_{lj}A_{ik}+\delta_{lk}A_{ij})-4f\right]\notag\\
 &+\left[D_lG_{lijk}+(2a-1)q_lG_{lijk}\right]X_{ijk}\notag\\
 &+w\left[D_lM_l+(2a+1)q_lM_l\right]
       +z\left[D_lP_l+2aq_lP_l\right]\notag\\
 &+f|\ell|^2-(D_lf+2aq_lf)\ell_l
 \ \ge\ \epsilon_*(z^2+|\ell|^2).\label{cert:interior}
\end{align}
In particular, $f\ge\epsilon_*$ whenever $g^2+|p|^2\le1$.

For every symmetric trace-free unit matrix $A$ and vector $p$ satisfying
\[
 g=0,\qquad |p|=1,\qquad Ap=\lambda p,\qquad\lambda<0,
\]
define
\[
 S_0=\lambda A-A^2,\qquad
 V_{ijk}=p_i(S_0)_{jk}+p_j(S_0)_{ki}+p_k(S_0)_{ij}+p_ip_jp_k.
\]
Then
\begin{equation}\label{cert:boundary}
 \mathcal S:=-\lambda f+U_{jk}(S_0+p\otimes p)_{jk}
                 +C_{lijk}p_lV_{ijk}+p_lM_l
 \ \ge\ \delta_*.
\end{equation}
Moreover, $U$ and $K$ are even in $p$, and $P=0$ when $g=0$.
\end{proposition}

We prove Proposition~\ref{cert:universal} in
Section~\ref{sec:certificate-proof}. First, Section~\ref{sec:flux}
derives its geometric consequences using only the stated properties.
The calculation explains the Laplacian and skew contractions in
$G$, the three coefficients of $q$, and the free vector $\ell$.
